\documentclass[../../main.tex]{subfiles}% 注意这里的文件路径不能用 ./main.tex ,否则用latexmk编译子文件会报错
\graphicspath{{\subfix{./image/}}{\subfix{../../image/}}} % 指定图片引用路径,后续可以直接使用图片文件名
% 如果图片存放在上级目录,则使用 ../../image/ 作为备用路径

% 例如:
% \begin{figure}[H]
% \centering
% \includegraphics[width=0.3\textwidth]{图.png}
% \caption{}
% \label{figure:图}
% \end{figure}
% 注意:上述\label{}一定要放在\caption{}之后,否则引用图片序号会只会显示??.

\begin{document}

\section{草稿四}

\begin{lemma}[片状函数的弱二阶导数]\label{lemma:ch05-piecewise-W2p}
设 $B=B(0,r)\subseteq\mathbb{R}^{n}$，$B^{+}\triangleq B\cap\{x_n>0\}$，$B^{-}\triangleq B\cap\{x_n<0\}$，$1\leqslant p\leqslant\infty$。若 $f\in C^{1}(B)$，且 $f|_{B^{\pm}}\in W^{2,p}(B^{\pm})$，则 $f\in W^{2,p}(B)$，并且对每个 $|\alpha|\leqslant2$，
\begin{align*}
D^{\alpha}f=
\begin{cases}
D^{\alpha}(f|_{B^{+}}), & \text{a.e. 于 }B^{+},\\
D^{\alpha}(f|_{B^{-}}), & \text{a.e. 于 }B^{-}.
\end{cases}
\end{align*}
\end{lemma}
\begin{proof}
当 $|\alpha|\leqslant1$ 时，由 $f\in C^{1}(B)$ 直接得到。下面设 $|\alpha|=2$，$\alpha=(i,j)$，取 $\varphi\in C_{c}^{\infty}(B)$。由于 $\varphi$ 在 $\partial B$ 上为零，在 $B^{+}$ 与 $B^{-}$ 上分别分部积分两次，得到
\begin{align*}
\int_{B^{\pm}}fD_iD_j\varphi\,\mathrm{d}x
=\int_{B^{\pm}}D_iD_jf\,\varphi\,\mathrm{d}x
+\int_{\Sigma}fD_j\varphi\,\nu_i^{\pm}\,\mathrm{d}S
-\int_{\Sigma}D_if\,\varphi\,\nu_j^{\pm}\,\mathrm{d}S,
\end{align*}
其中 $\Sigma\triangleq B\cap\{x_n=0\}$，$\nu^{\pm}$ 是 $\partial B^{\pm}$ 的单位外法向量。由于 $\nu^{-}=-\nu^{+}$，且 $f$ 与 $D_if$ 在 $\Sigma$ 上连续，右端沿 $\Sigma$ 的边界项相消。于是
\begin{align*}
\int_{B}fD_iD_j\varphi\,\mathrm{d}x
=\int_{B^{+}}D_iD_jf\,\varphi\,\mathrm{d}x
+\int_{B^{-}}D_iD_jf\,\varphi\,\mathrm{d}x.
\end{align*}
由 $\varphi$ 的任意性，弱导数 $D_iD_jf$ 存在，且几乎处处等于逐块二阶导数。证毕。
\end{proof}

\begin{theorem}[$W^{2,p}$ 扩张定理]\label{theorem:ch05-W2p扩张}
设 $1\leqslant p\leqslant\infty$，$U\subseteq\mathbb{R}^{n}$ 是有界开集，$\partial U$ 为 $C^{2}$，并取有界开集 $V$ 使 $U\subset\subset V$。则 \cref{theorem:ch05-扩张定理} 中构造的扩张算子 $E$ 限制在 $W^{2,p}(U)$ 上仍为有界线性算子，即
\begin{align*}
E\bigl(W^{2,p}(U)\bigr)\subseteq W^{2,p}(\mathbb{R}^{n}),
\end{align*}
且存在常数 $C=C(n,p,U,V)>0$，使得对每个 $u\in W^{2,p}(U)$，
\begin{align}
\|Eu\|_{W^{2,p}(\mathbb{R}^{n})}\leqslant C\|u\|_{W^{2,p}(U)}.
\label{eq:ch05-W2p-extension-estimate}
\end{align}
此外，$Eu=u$ a.e. 于 $U$，且 $\operatorname{spt}Eu\subset\subset V$。
\end{theorem}
\begin{proof}
\noindent$\mathbf{Step}\bm{1}$.{\heiti 证明半空间中的二阶反射估计.}

设 $B=B(0,r)$，$B^{+}=B\cap\{x_n>0\}$，$B^{-}=B\cap\{x_n<0\}$，并记 $x'=(x_1,\cdots,x_{n-1})$。对 $u\in C^{\infty}(\overline{B^{+}})$，按 \eqref{align:ch05-反射延拓} 定义
\begin{align*}
\overline u(x)\triangleq
\begin{cases}
u(x), & x\in B^{+},\\
-3u(x',-x_n)+4u\left(x',-\dfrac{x_n}{2}\right), & x\in B^{-}.
\end{cases}
\end{align*}
在 $\{x_n=0\}$ 上两式相等，故定义无歧义。直接求导并令 $x_n\to0^{\pm}$ 可知 $\overline u$ 及其一阶偏导数在 $\{x_n=0\}$ 两侧有相同的连续极限；再由中值定理（法向）和边界上的 $C^{1}$ 性（切向），$\overline u\in C^{1}(B)$，且 $\overline u|_{B^{\pm}}\in C^{\infty}(B^{\pm})$。由 \cref{lemma:ch05-piecewise-W2p}，$\overline u\in W^{2,p}(B)$。

对 $x\in B^{-}$ 与 $|\alpha|=2$，直接求导得
\begin{align*}
D^{\alpha}\overline u(x)=
\begin{cases}
-3D^{\alpha}u(x',-x_n)+4D^{\alpha}u\left(x',-\dfrac{x_n}{2}\right), & \alpha=(i,j),\ 1\leqslant i,j\leqslant n-1,\\[4pt]
3D^{\alpha}u(x',-x_n)-2D^{\alpha}u\left(x',-\dfrac{x_n}{2}\right), & \alpha=(i,n)\text{ 或 }(n,i),\ 1\leqslant i\leqslant n-1,\\[4pt]
-3D^{\alpha}u(x',-x_n)+D^{\alpha}u\left(x',-\dfrac{x_n}{2}\right), & \alpha=(n,n).
\end{cases}
\end{align*}

当 $1\leqslant p<\infty$ 时，由 $|a+b|^p\leqslant 2^{p-1}(|a|^p+|b|^p)$ 得
\begin{align*}
\int_{B^{-}}|D^{\alpha}\overline u|^p\,\mathrm{d}x
\leqslant C_p\left(3^p\int_{B^{-}}|D^{\alpha}u(x',-x_n)|^p\,\mathrm{d}x
+4^p\int_{B^{-}}\left|D^{\alpha}u\left(x',-\frac{x_n}{2}\right)\right|^p\,\mathrm{d}x\right),
\end{align*}
其中 $C_p=2^{p-1}$。对右端第一个积分作变量替换 $y_n=-x_n$，对第二个积分作变量替换 $y_n=-x_n/2$，并注意第二个变换的 Jacobi 行列式绝对值为 $2$，得到
\begin{align*}
\int_{B^{-}}|D^{\alpha}u(x',-x_n)|^p\,\mathrm{d}x
&=\int_{B^{+}}|D^{\alpha}u(y)|^p\,\mathrm{d}y,\\
\int_{B^{-}}\left|D^{\alpha}u\left(x',-\frac{x_n}{2}\right)\right|^p\,\mathrm{d}x
&=2\int_{E}|D^{\alpha}u(y)|^p\,\mathrm{d}y
\leqslant2\int_{B^{+}}|D^{\alpha}u(y)|^p\,\mathrm{d}y,
\end{align*}
其中 $E=\{y\in\mathbb{R}^{n}:y_1^2+\cdots+y_{n-1}^2+4y_n^2<r^2,\ y_n>0\}\subseteq B^{+}$。因此
\begin{align*}
\int_{B^{-}}|D^{\alpha}\overline u|^p\,\mathrm{d}x
\leqslant C_p(3^p+2\cdot4^p)\int_{B^{+}}|D^{\alpha}u|^p\,\mathrm{d}x.
\end{align*}

当 $p=\infty$ 时，同样有
\begin{align*}
\|D^{\alpha}\overline u\|_{L^{\infty}(B)}
\leqslant C\|D^{\alpha}u\|_{L^{\infty}(B^{+})},\quad |\alpha|\leqslant2.
\end{align*}
为避免常数记号，$p=\infty$ 时以下使用等价范数
\begin{align*}
\|v\|'_{W^{2,\infty}}\triangleq\max_{|\alpha|\leqslant2}\|D^{\alpha}v\|_{L^{\infty}};
\end{align*}
它与 \cref{definition:ch05-Wkp范数} 中的范数等价，故结论不变。

对 $|\alpha|\leqslant1$，同样的计算给出
\begin{align*}
\|D^{\alpha}\overline u\|_{L^{p}(B)}
\leqslant C\|D^{\alpha}u\|_{L^{p}(B^{+})}.
\end{align*}
合并 $|\alpha|\leqslant2$ 的估计，得到
\begin{align}
\|\overline u\|_{W^{2,p}(B)}
\leqslant C\|u\|_{W^{2,p}(B^{+})}.
\label{eq:ch05-W2p-local-reflection}
\end{align}

\noindent$\mathbf{Step}\bm{2}$.{\heiti 拉直边界并证明局部估计.}

固定 $x_0\in\partial U$。因为 $\partial U$ 为 $C^{2}$，在 \cref{definition:appC-Ck边界} 中可取 $\gamma\in C^{2}(\mathbb{R}^{n-1})$，从而 \cref{theorem:ch05-扩张定理} 步骤5中的 $\Phi_1,\Psi_1$ 均为 $C^{2}$ 微分同胚。重复该步骤的构造，得到 $x_0$ 的开邻域 $W\subset\subset V$、开球 $B=B(0,\rho)$ 以及 $C^{2}$ 微分同胚 $\Phi:W\to B$，其逆为 $\Psi:B\to W$，使得
\begin{align*}
\Phi(W)=B,\qquad \Phi(W\cap U)=B^{+},\qquad \Phi(W\cap\partial U)=B\cap\{y_n=0\}.
\end{align*}

设 $u\in C^{\infty}(\overline U)$，令 $u'=u\circ\Psi$，并按步骤1的反射公式定义 $\overline{u'}$；再令 $\overline u=\overline{u'}\circ\Phi$。由步骤1，
\begin{align*}
\|\overline{u'}\|_{W^{2,p}(B)}
\leqslant C\|u'\|_{W^{2,p}(B^{+})}.
\end{align*}
对 $v\in W^{2,p}(B)$，链式法则为
\begin{align}
D_{ij}(v\circ\Phi)
=\sum_{k,l=1}^{n}(D_{kl}v)\circ\Phi\,D_i\Phi_kD_j\Phi_l
+\sum_{k=1}^{n}(D_kv)\circ\Phi\,D_{ij}\Phi_k.
\label{eq:ch05-W2p-chain}
\end{align}
因为 $\Phi,\Psi$ 在相应的闭包上为 $C^{2}$，且 $|\det D\Phi|,|\det D\Psi|$ 有正的上、下界，由 \eqref{eq:ch05-W2p-chain} 与变量替换，存在 $C_{\Phi},C_{\Psi}>0$，使得
\begin{align*}
\|v\circ\Phi\|_{W^{2,p}(W)}
&\leqslant C_{\Phi}\|v\|_{W^{2,p}(B)},\\
\|w\circ\Psi\|_{W^{2,p}(B^{+})}
&\leqslant C_{\Psi}\|w\|_{W^{2,p}(W\cap U)}.
\end{align*}
$p=\infty$ 时同样成立。于是
\begin{align*}
\|\overline u\|_{W^{2,p}(W)}
&\leqslant C_{\Phi}\|\overline{u'}\|_{W^{2,p}(B)}
\leqslant C_{\Phi}C\|u'\|_{W^{2,p}(B^{+})}
\leqslant C_{\Phi}CC_{\Psi}\|u\|_{W^{2,p}(W\cap U)}.
\end{align*}
记 $C'=C_{\Phi}CC_{\Psi}$，则
\begin{align}
\|\overline u\|_{W^{2,p}(W)}
\leqslant C'\|u\|_{W^{2,p}(W\cap U)}.
\label{eq:ch05-W2p-local-estimate}
\end{align}
若 $x\in W\cap U$，则 $\Phi(x)\in B^{+}$，于是
\begin{align*}
\overline u(x)=\overline{u'}(\Phi(x))=u'(\Phi(x))=u(\Psi(\Phi(x)))=u(x).
\end{align*}

\noindent$\mathbf{Step}\bm{3}$.{\heiti 用单位分解拼接局部延拓.}

由 $\partial U$ 紧，取有限个点 $x_i^0\in\partial U$ 及相应的开集 $W_i\subset\subset V$（$i=1,\cdots,N$）覆盖 $\partial U$。再取 $W_0\subset\subset U$ 使 $U\subseteq\cup_{i=0}^{N}W_i$。取从属于 $\{W_i\}_{i=0}^{N}$ 的光滑单位分解 $\{\zeta_i\}_{i=0}^{N}$，满足 $\zeta_i\in C_c^{\infty}(W_i)$ 且 $\sum_{i=0}^{N}\zeta_i=1$ 于 $U$。对 $u\in C^{\infty}(\overline U)$，令 $\overline u_0=u$，$\overline u_i$（$i\geqslant1$）为步骤2得到的局部延拓，并定义
\begin{align*}
E_0u\triangleq\sum_{i=0}^{N}\zeta_i\overline u_i.
\end{align*}
由 Leibniz 公式，对 $|\alpha|\leqslant2$，
\begin{align*}
D^{\alpha}(\zeta_i\overline u_i)
=\sum_{\beta\leqslant\alpha}\binom{\alpha}{\beta}D^{\beta}\zeta_i\,D^{\alpha-\beta}\overline u_i.
\end{align*}
由于 $\zeta_i\in C_c^{\infty}(W_i)$，故存在 $C_i>0$，使得
\begin{align*}
\|\zeta_i\overline u_i\|_{W^{2,p}(W_i)}
\leqslant C_i\|\overline u_i\|_{W^{2,p}(W_i)}.
\end{align*}
于是由 \eqref{eq:ch05-W2p-local-estimate}，
\begin{align}
\|E_0u\|_{W^{2,p}(\mathbb{R}^{n})}
&\leqslant\sum_{i=0}^{N}\|\zeta_i\overline u_i\|_{W^{2,p}(W_i)}
\leqslant C\sum_{i=0}^{N}\|\overline u_i\|_{W^{2,p}(W_i)}
\leqslant C\|u\|_{W^{2,p}(U)},\quad u\in C^{\infty}(\overline U).
\label{eq:ch05-W2p-smooth-estimate}
\end{align}
并且
\begin{align*}
\operatorname{spt}E_0u
\subseteq\bigcup_{i=0}^{N}\operatorname{spt}\zeta_i
\subseteq\bigcup_{i=0}^{N}W_i
\subset\subset V,
\end{align*}
而
\begin{align*}
E_0u(x)=\sum_{i=0}^{N}\zeta_i(x)u(x)=u(x),\quad x\in U.
\end{align*}

\noindent$\mathbf{Step}\bm{4}$.{\heiti 逼近延拓到 $W^{2,p}(U)$，$1\leqslant p<\infty$.}

设 $1\leqslant p<\infty$，$u\in W^{2,p}(U)$。由 \cref{theorem:ch05-直到边界光滑的函数整体逼近定理}，取 $u_m\in C^{\infty}(\overline U)$ 使 $u_m\to u$ 于 $W^{2,p}(U)$。由 \eqref{eq:ch05-W2p-smooth-estimate} 与线性性，
\begin{align*}
\|E_0u_m-E_0u_l\|_{W^{2,p}(\mathbb{R}^{n})}
\leqslant C\|u_m-u_l\|_{W^{2,p}(U)}.
\end{align*}
故 $\{E_0u_m\}$ 是 $W^{2,p}(\mathbb{R}^{n})$ 中的 Cauchy 列。由 \cref{theorem:ch05-WkpBanach}，存在 $v\in W^{2,p}(\mathbb{R}^{n})$ 使 $E_0u_m\to v$。

另一方面，对光滑函数 $E_0=E$；而 $E$ 在 $W^{1,p}(U)$ 上有界，故
\begin{align*}
E_0u_m=Eu_m\to Eu\quad\text{于 }W^{1,p}(\mathbb{R}^{n}).
\end{align*}
于是 $v=Eu$ a.e.，从而 $Eu\in W^{2,p}(\mathbb{R}^{n})$，且
\begin{align*}
\|Eu\|_{W^{2,p}(\mathbb{R}^{n})}
=\lim_{m\to\infty}\|E_0u_m\|_{W^{2,p}(\mathbb{R}^{n})}
\leqslant C\|u\|_{W^{2,p}(U)}.
\end{align*}
良定义性由 $W^{2,p}$ 极限的唯一性和 $E$ 在 $W^{1,p}$ 上已良定义得到。

\noindent$\mathbf{Step}\bm{5}$.{\heiti 处理 $p=\infty$.}

设 $u\in W^{2,\infty}(U)$。由 Sobolev 嵌入 $W^{2,\infty}(U)\subseteq C^{1}(\overline U)$，取 $u$ 的 $C^{1}$ 代表元，仍记为 $u$，且 $D^{2}u\in L^{\infty}(U)$。对步骤2的每个局部图，令 $u'=u\circ\Psi$。则 $u'\in C^{1}(\overline{B^{+}})$，且 $D^{2}u'\in L^{\infty}(B^{+})$。按步骤1的反射公式定义 $\overline{u'}$。同步骤1，$\overline{u'}\in C^{1}(B)$，并由 \cref{lemma:ch05-piecewise-W2p} 得 $\overline{u'}\in W^{2,\infty}(B)$，且
\begin{align*}
\|\overline{u'}\|_{W^{2,\infty}(B)}
\leqslant C\|u'\|_{W^{2,\infty}(B^{+})}.
\end{align*}
令 $\overline u=\overline{u'}\circ\Phi$。由步骤2的链式法则（$p=\infty$）得 $\overline u\in W^{2,\infty}(W)$，且
\begin{align*}
\|\overline u\|_{W^{2,\infty}(W)}
\leqslant C\|u\|_{W^{2,\infty}(W\cap U)}.
\end{align*}
重复步骤3的有限覆盖与单位分解，得到 $E_0u\in W^{2,\infty}(\mathbb{R}^{n})$，满足
\begin{align*}
\|E_0u\|_{W^{2,\infty}(\mathbb{R}^{n})}
\leqslant C\|u\|_{W^{2,\infty}(U)}.
\end{align*}
由于 $p=\infty$ 时 \cref{theorem:ch05-扩张定理} 中的 $E$ 正是由同一局部反射公式和单位分解定义的，故 $Eu=E_0u$。因此
\begin{align*}
\|Eu\|_{W^{2,\infty}(\mathbb{R}^{n})}
\leqslant C\|u\|_{W^{2,\infty}(U)}.
\end{align*}
综上，对每个 $1\leqslant p\leqslant\infty$，定理得证。
\end{proof}

\begin{remark}\label{remark:ch05-延拓注}
\leavevmode
\begin{enumerate}[(i)]
\item \cref{theorem:ch05-W2p扩张} 表明，当 $\partial U$ 为 $C^{2}$ 时，\cref{theorem:ch05-扩张定理} 中的扩张算子 $E$ 还把 $W^{2,p}(U)$ 有界地映入 $W^{2,p}(\mathbb{R}^{n})$。证明的关键是：反射公式 \eqref{align:ch05-反射延拓} 只保证 $\overline u\in C^{1}$，但其弱二阶导数在 $\{x_n=0\}$ 两侧分别存在；由于一阶偏导数穿过该平面连续，分部积分时边界项相消，不会出现奇异测度。
\item 上述构造并未为 $k>2$ 的 Sobolev 空间 $W^{k,p}(U)$ 提供延拓。这需要更复杂的高阶反射技巧。
\end{enumerate}
\end{remark}










\end{document}